% Shared layout for every language version of the paper.
% Fonts are loaded from paper/fonts/, so every machine and CI build sets the same glyphs.
\usepackage{fontspec}
\usepackage[english,russian]{babel}
\babelfont{rm}[Path=fonts/ptserif/, Extension=.ttf, UprightFont=*-Regular, BoldFont=*-Bold,
  ItalicFont=*-Italic, BoldItalicFont=*-BoldItalic]{PT_Serif-Web}
\babelfont{sf}[Path=fonts/ptsans/, Extension=.ttf, UprightFont=*-Regular, BoldFont=*-Bold,
  ItalicFont=*-Italic, BoldItalicFont=*-BoldItalic]{PT_Sans-Web}
\babelfont{tt}[Path=fonts/ptmono/, Extension=.ttf, Scale=0.88]{PTM55FT}
\usepackage[a4paper,top=20mm,bottom=22mm,left=20mm,right=20mm]{geometry}
\usepackage{amsmath}
\usepackage{booktabs}
\usepackage{array}
\usepackage{enumitem}
\usepackage{titlesec}
\usepackage{caption}
\usepackage{tikz}
\usetikzlibrary{arrows.meta}
\usepackage{xcolor}
\usepackage[hidelinks]{hyperref}
\usepackage{microtype}

\setlength{\parindent}{0pt}
\setlength{\parskip}{3pt plus 1pt}
\setlength{\emergencystretch}{4em}

\titleformat{\section}{\normalfont\fontsize{11.5}{14}\bfseries}{\thesection.}{0.5em}{}
\titlespacing*{\section}{0pt}{8pt plus 2pt minus 2pt}{3pt}
\titleformat{\subsection}{\normalfont\fontsize{10}{12}\bfseries}{\thesubsection.}{0.5em}{}
\titlespacing*{\subsection}{0pt}{5pt plus 2pt minus 1pt}{2pt}

\setlist[enumerate]{leftmargin=2.0em,itemsep=0.5pt,topsep=2pt,parsep=0pt}
\setlist[itemize]{leftmargin=1.5em,itemsep=0.5pt,topsep=2pt,parsep=0pt}
\captionsetup{font=small,labelfont=bf,labelsep=period,skip=3pt}
\newcolumntype{L}[1]{>{\raggedright\arraybackslash}p{#1}}
\setlength{\textfloatsep}{7pt plus 2pt minus 2pt}
\setlength{\floatsep}{8pt plus 2pt minus 2pt}

\pagestyle{plain}

\tikzset{
  box/.style={draw,rectangle,thick,align=center,inner sep=2.5pt,font=\footnotesize,minimum width=27mm,minimum height=8mm},
  zone/.style={draw,dashed,rounded corners=2pt},
  zl/.style={font=\scriptsize\itshape,anchor=north west,inner sep=2pt},
  lbl/.style={font=\scriptsize,align=center,fill=white,inner sep=1pt},
  arr/.style={-{Latex[length=1.8mm]},thick},
}

% Version line from version.tex, which `make pdf` writes from git (.scripts/version.sh):
% a release tag gives "Октябрь 2026, версия 0.1.0"; any other build is marked as a draft.
\newcommand{\papermonth}{\ifcase\month\or Январь\or Февраль\or Март\or Апрель\or Май\or Июнь\or
  Июль\or Август\or Сентябрь\or Октябрь\or Ноябрь\or Декабрь\fi}
\newcommand{\paperversionline}{\papermonth\ \the\year,
  \ifpaperrelease версия \paperversion\else черновик \paperversion\fi}
